\section{Physical network model}
\label{sec:network}

The network case contains 11 buses, 6 branches, 4 transformers, 9 generators plus the
POI swing source, and 4 loads. All equipment is assigned to area~26 (LADWP) and
zone~283 (LA~IPP). The system base is 100\,MV$\cdot$A. Figure~\ref{fig:oneline}
shows the one-line arrangement.

\begin{figure}[!ht]
\centering
\resizebox{\linewidth}{!}{%
\begin{tikzpicture}[x=1cm,y=1cm,font=\footnotesize]
  % POI
  \draw[bus] (-1.2,10) -- (1.2,10) node[right,xshift=2mm]{\textbf{1 POI} 345\,kV\quad(swing, $U_{sched}=1.0493$)};
  \draw (0,10) -- (0,8.6);
  \node[right,align=left] at (0.15,9.3) {3.0\,mi, 1590 kcmil Lapwing\\1077.7\,MV$\cdot$A};
  % GEN-TIE busbar
  \draw[bus] (-6.4,8.6) -- (6.4,8.6) node[right,xshift=2mm]{\textbf{18 GEN-TIE} 345\,kV};
  % tie to Phase 1
  \draw (-6.0,8.6) -- (-6.0,7.6) -- (-7.4,7.6);
  \draw[bus] (-7.4,7.1) -- (-7.4,8.1);
  \node[left,align=right] at (-7.5,7.6) {\textbf{2 HV-MPT-P1}\\345\,kV (Phase~1)};
  \node[above,align=center] at (-6.7,7.65) {3.2\,mi\\Lapwing};
  % four feeders
  \foreach \x/\b/\g/\len/\gsu/\u in {
      -4.5/19/17/0.60/HV-GSU-P2-1/{4, 5},
      -1.5/20/22/0.74/HV-GSU-P2-2/{6, 7},
       1.5/21/23/2.81/HV-GSU-P2-3/{8, 9},
       4.5/28/25/2.88/HV-GSU-P2-4/{10, 11}}{
    \draw (\x,8.6) -- (\x,7.0);
    \node[right,align=left] at (\x+0.08,7.85) {\len\,mi\\Waxwing};
    \draw[bus] (\x-0.8,7.0) -- (\x+0.8,7.0);
    \node[above left,align=right,yshift=-0.3mm] at (\x-0.8,7.0) {\textbf{\b}};
    \node[below,align=center] at (\x,6.95) {\gsu\\345\,kV};
    % transformer
    \draw (\x,6.35) -- (\x,5.95);
    \draw (\x,5.6) circle (0.35);
    \draw (\x,5.1) circle (0.35);
    \draw (\x,4.75) -- (\x,4.4);
    \node[right,align=left,font=\scriptsize] at (\x+0.4,5.35) {150\,MV$\cdot$A\\$X_T = 14.5\,\%$\\362/13.8\,kV};
    % gen bus
    \draw[bus] (\x-0.9,4.4) -- (\x+0.9,4.4);
    \node[above left,align=right,yshift=-0.3mm] at (\x-0.9,4.4) {\textbf{\g}};
    \node[below left,align=right,font=\scriptsize,yshift=0.3mm] at (\x-0.9,4.4) {13.8\,kV};
    % two generators
    \draw (\x-0.55,4.4) -- (\x-0.55,3.75);
    \draw (\x-0.55,3.35) circle (0.4) node{G};
    \draw (\x+0.55,4.4) -- (\x+0.55,3.75);
    \draw (\x+0.55,3.35) circle (0.4) node{G};
    % load arrow
    \draw[-{Latex[length=2mm]}] (\x,4.4) -- (\x,2.6);
    \node[below,align=center] at (\x,2.55) {aux\\load};
    \node[below,align=center] at (\x,1.9) {units \u};
  }
\end{tikzpicture}
}
\caption{One-line diagram of the Phase~2 network case. Bus numbers are those of the \texttt{.epc} file; line lengths and conductor types are taken from the branch records.}
\label{fig:oneline}
\end{figure}

\subsection{Buses}

Table~\ref{tab:bus} lists the buses. Bus~1 is the swing bus (PSLF type~0) with a
scheduled voltage of 1.0493\,pu; the four generator buses are voltage-controlled
(type~2) at 1.00\,pu, regulating their own terminals; the remaining 345\,kV buses are
load buses (type~1). Every bus carries the same 0.95--1.05\,pu voltage band.

\begin{table}[!ht]
\centering
\begin{tabularx}{\linewidth}{@{}r l r c r X@{}}
\toprule
\multicolumn{1}{c}{\textbf{No.}} & \textbf{Name} & \multicolumn{1}{c}{\textbf{kV}} & \multicolumn{1}{c}{\textbf{Type}} & \multicolumn{1}{c}{\textbf{$U_{sched}$}} & \textbf{Function} \\
\midrule
 1 & POI         & 345.0 & 0 & 1.0493 & Point of interconnection, swing bus \\
 2 & HV-MPT-P1   & 345.0 & 1 & --     & Phase~1 main power transformer HV bus \\
18 & GEN-TIE     & 345.0 & 1 & --     & Phase~2 345\,kV collector bus \\
19 & HV-GSU-P2-1 & 345.0 & 1 & --     & GSU~1 HV terminal \\
20 & HV-GSU-P2-2 & 345.0 & 1 & --     & GSU~2 HV terminal \\
21 & HV-GSU-P2-3 & 345.0 & 1 & --     & GSU~3 HV terminal \\
28 & HV-GSU-P2-4 & 345.0 & 1 & --     & GSU~4 HV terminal \\
17 & GEN-P2-U4-5 & 13.8  & 2 & 1.0000 & Generator bus, units 4 and 5, aux load \\
22 & GEN-P2-U6-7 & 13.8  & 2 & 1.0000 & Generator bus, units 6 and 7, aux load \\
23 & GEN-P2-U8-9 & 13.8  & 2 & 1.0000 & Generator bus, units 8 and 9, aux load \\
25 & GEN-P2-U10-1 & 13.8 & 2 & 1.0000 & Generator bus, units 10 and 11, aux load \\
\bottomrule
\end{tabularx}
\caption{Bus data (all buses in area 26 LADWP, zone 283 LA IPP).}
\label{tab:bus}
\end{table}

\subsection{Branches}

The six 345\,kV branches are given in Table~\ref{tab:branch}. Impedances are in per
unit on the 100\,MV$\cdot$A system base; \texttt{charge} is the total line charging
$B$. The POI and Phase~1 ties use 1590\,kcmil Lapwing ACSR; the four GSU feeders use
266.8\,kcmil Waxwing ACSR. The GEN-TIE to HV-MPT-P1 branch carries a metering flag
at the GEN-TIE end and a 90 line-section identifier in the record, which is how the
Phase~1/Phase~2 boundary is marked in the case.

\begin{table}[!ht]
\centering
\footnotesize\setlength{\tabcolsep}{4pt}\begin{tabularx}{\linewidth}{@{}l l r r r r r r r r@{}}
\toprule
\textbf{From} & \textbf{To} & \multicolumn{1}{c}{\textbf{$R$}} & \multicolumn{1}{c}{\textbf{$X$}} & \multicolumn{1}{c}{\textbf{$B$}} & \multicolumn{1}{c}{\textbf{Rate~1}} & \multicolumn{1}{c}{\textbf{Rate~2}} & \multicolumn{1}{c}{\textbf{Rate~3}} & \multicolumn{1}{c}{\textbf{Rate~4}} & \multicolumn{1}{c}{\textbf{mi}} \\
\midrule
 1 POI     & 18 GEN-TIE     & 0.000145 & 0.001841 & 0.021185 & 1077.7 & 1291.2 & 1351.8 & 1408.4 & 3.00 \\
18 GEN-TIE &  2 HV-MPT-P1   & 0.000155 & 0.001963 & 0.022597 & 1077.7 & 1291.2 & 1351.8 & 1408.4 & 3.20 \\
18 GEN-TIE & 19 HV-GSU-P2-1 & 0.000159 & 0.000416 & 0.003687 &  144.6 &  190.9 &  204.0 &  216.2 & 0.60 \\
18 GEN-TIE & 20 HV-GSU-P2-2 & 0.000196 & 0.000513 & 0.004548 &  144.6 &  190.9 &  204.0 &  216.2 & 0.74 \\
18 GEN-TIE & 21 HV-GSU-P2-3 & 0.000745 & 0.001949 & 0.017269 &  144.6 &  190.9 &  204.0 &  216.2 & 2.81 \\
18 GEN-TIE & 28 HV-GSU-P2-4 & 0.000763 & 0.001997 & 0.017699 &  144.6 &  190.9 &  204.0 &  216.2 & 2.88 \\
\bottomrule
\end{tabularx}
\caption{Branch data (pu on 100\,MV$\cdot$A; ratings in MV$\cdot$A; the two 3-mile ties are Lapwing, the four feeders Waxwing).}
\label{tab:branch}
\end{table}

\subsection{Transformers}

All four GSUs are identical two-winding units, Table~\ref{tab:xfmr}. The impedance is
entered on the transformer base of 150\,MV$\cdot$A, so on the 100\,MV$\cdot$A system
base it becomes
\begin{equation}
  X_{T,sys} = 0.14496 \cdot \frac{100}{150} = 0.09664\,\mathrm{pu},
  \qquad
  R_{T,sys} = 0.00345 \cdot \frac{100}{150} = 0.00230\,\mathrm{pu}.
\end{equation}
The HV winding tap is 362\,kV against a 345\,kV bus base, i.e.\ a fixed
$362/345 = 1.0493$\,pu tap position, which matches the scheduled POI voltage. The
transformers are of PSLF type~1 (fixed tap, no regulation) with a 200\,MV$\cdot$A
rating in both rate slots.

\begin{table}[!ht]
\centering
\small\begin{tabularx}{\linewidth}{@{}l l c r r r r r@{}}
\toprule
\textbf{HV bus} & \textbf{LV bus} & \multicolumn{1}{c}{\textbf{ck}} & \multicolumn{1}{c}{\textbf{$S_{base}$}} & \multicolumn{1}{c}{\textbf{$R_T$}} & \multicolumn{1}{c}{\textbf{$X_T$}} & \multicolumn{1}{c}{\textbf{Tap HV/LV}} & \multicolumn{1}{c}{\textbf{Rating}} \\
 & & & \multicolumn{1}{c}{MV$\cdot$A} & \multicolumn{1}{c}{pu} & \multicolumn{1}{c}{pu} & \multicolumn{1}{c}{kV} & \multicolumn{1}{c}{MV$\cdot$A} \\
\midrule
19 HV-GSU-P2-1 & 17 GEN-P2-U4-5  & 1  & 150 & 0.00345 & 0.14496 & 362.0 / 13.8 & 200 \\
20 HV-GSU-P2-2 & 22 GEN-P2-U6-7  & 1  & 150 & 0.00345 & 0.14496 & 362.0 / 13.8 & 200 \\
21 HV-GSU-P2-3 & 23 GEN-P2-U8-9  & 11 & 150 & 0.00345 & 0.14496 & 362.0 / 13.8 & 200 \\
28 HV-GSU-P2-4 & 25 GEN-P2-U10-1 & 1  & 150 & 0.00345 & 0.14496 & 362.0 / 13.8 & 200 \\
\bottomrule
\end{tabularx}
\caption{Generator step-up transformer data.}
\label{tab:xfmr}
\end{table}

\subsection{Generators}

Table~\ref{tab:gen} gives the generator dispatch and limits. Each unit regulates
its own terminal bus at 1.00\,pu with a full share of the reactive requirement
(\texttt{prf} $=$ \texttt{qrf} $= 1.0$). The dynamic source impedance
\texttt{gen\_x} is 1.0\,pu on the machine base for every unit and is overridden by
the machine model at initialisation. The POI source is a 99\,999\,MV$\cdot$A swing
generator with $\pm$9999\,MW and $\pm$9999\,Mvar limits, i.e.\ an ideal source.

\begin{table}[!ht]
\centering
\begin{tabularx}{\linewidth}{@{}r c X r r r r r r@{}}
\toprule
\multicolumn{1}{c}{\textbf{Bus}} & \multicolumn{1}{c}{\textbf{id}} & \textbf{Long name} & \multicolumn{1}{c}{\textbf{$P_{gen}$}} & \multicolumn{1}{c}{\textbf{$P_{max}$}} & \multicolumn{1}{c}{\textbf{$P_{min}$}} & \multicolumn{1}{c}{\textbf{$Q_{max}$}} & \multicolumn{1}{c}{\textbf{$Q_{min}$}} & \multicolumn{1}{c}{\textbf{$S_{base}$}} \\
\midrule
17 & 1 & Cape Phase 2 -- Turboden -- Gen 4  & 56.659 & 65.000 & 0 & 45.118 & $-36.706$ & 76.471 \\
17 & 2 & Cape Phase 2 -- Turboden -- Gen 5  & 56.659 & 65.000 & 0 & 45.118 & $-36.706$ & 76.471 \\
22 & 1 & Cape Phase 2 -- Turboden -- Gen 6  & 56.659 & 65.000 & 0 & 45.118 & $-36.706$ & 76.471 \\
22 & 2 & Cape Phase 2 -- Baker H -- Gen 7   & 55.402 & 60.461 & 0 & 46.235 & $-38.410$ & 71.130 \\
23 & 1 & Cape Phase 2 -- Baker H -- Gen 8   & 55.402 & 60.461 & 0 & 46.235 & $-38.410$ & 71.130 \\
23 & 2 & Cape Phase 2 -- Baker H -- Gen 9   & 55.402 & 60.461 & 0 & 46.235 & $-38.410$ & 71.130 \\
25 & 1 & Cape Phase 2 -- Baker H -- Gen 10  & 55.402 & 60.461 & 0 & 46.235 & $-38.410$ & 71.130 \\
25 & 2 & Cape Phase 2 -- Baker H -- Gen 11  & 55.402 & 60.461 & 0 & 46.235 & $-38.410$ & 71.130 \\
\midrule
 1 & 1 & POI (swing source) & 0 & 9999 & $-9999$ & 9999 & $-9999$ & 99999 \\
\bottomrule
\end{tabularx}
\caption{Generator data (MW, Mvar and MV$\cdot$A).}
\label{tab:gen}
\end{table}

\subsection{Loads}

Each generator bus carries a constant-power auxiliary load representing the plant's
own consumption (geothermal brine pumping and balance of plant), Table~\ref{tab:load}.
The loads are modelled entirely as constant $P$ and $Q$; the current- and
impedance-dependent components are zero.

\begin{table}[!ht]
\centering
\begin{tabularx}{\linewidth}{@{}r l X r r@{}}
\toprule
\multicolumn{1}{c}{\textbf{Bus}} & \textbf{Name} & \textbf{Long name} & \multicolumn{1}{c}{\textbf{$P$ (MW)}} & \multicolumn{1}{c}{\textbf{$Q$ (Mvar)}} \\
\midrule
17 & GEN-P2-U4-5  & gen aux load & 35.60 & 17.24 \\
22 & GEN-P2-U6-7  & gen aux load & 34.36 & 16.64 \\
23 & GEN-P2-U8-9  & gen aux load & 33.12 & 16.04 \\
25 & GEN-P2-U10-1 & gen aux load & 33.12 & 16.04 \\
\midrule
\multicolumn{3}{l}{Total} & 136.20 & 65.96 \\
\bottomrule
\end{tabularx}
\caption{Auxiliary loads (constant power).}
\label{tab:load}
\end{table}

\subsection{Dispatch summary}

The scheduled dispatch of the nine units sums to
\begin{equation}
  P_{gross} = 3 \cdot 56.659\,\mathrm{MW} + 5 \cdot 55.402\,\mathrm{MW} = 446.98\,\mathrm{MW},
\end{equation}
against a combined capability of $3 \cdot 65.0 + 5 \cdot 60.461 = 497.3$\,MW. After
the 136.20\,MW of auxiliary load the net injection into the four GSUs is
$446.98 - 136.20 = 310.78$\,MW before transformer and line losses. The reactive
capability of the fleet spans $+366.5$ to $-302.2$\,Mvar at the terminals.
